\documentclass[11pt]{article}
\usepackage[T1]{fontenc}
\usepackage[margin=0.88in]{geometry}
\usepackage{amsmath,amssymb}
\usepackage{booktabs,tabularx}
\usepackage[hidelinks]{hyperref}
\usepackage{microtype}

\newcommand{\C}{\mathbb C}
\newcommand{\A}{\mathcal A}
\newcommand{\cS}{\mathcal S}

\title{The nonmetrizable finite-index expectation question}
\author{Literature-already-answered note for arXiv:2409.03531}
\date{11 August 2026}

\begin{document}
\maketitle

\paragraph{Status.}
\textbf{Answered negatively in the direct follow-up arXiv:2409.17807.}
The source paper proves existence over compact metrizable bases; the follow-up
constructs a compact nonmetrizable base and a continuous unital
subhomogeneous commutative bundle for which there is no conditional
expectation at all.  Thus the unrestricted compact-Hausdorff question has a
negative answer, while the metrizable theorem remains sharp in a meaningful
selection-theoretic sense.

\section*{The question and the residue left by the source paper}
Blanchard--Gogi\'{c} Problem~3.11, restated in \cite{source}, asks whether a
continuous unital subhomogeneous $C^*$-bundle
\(
  \A\twoheadrightarrow X
\)
over compact Hausdorff $X$ always makes the central inclusion
\[
 C(X)\hookrightarrow \Gamma(\A)
\]
admit a finite-index conditional expectation.  Here finite index means that
for an expectation $E$ there is a finite $K$ with $K E-\mathrm{id}\ge0$.
It also asks whether one can attain the optimal constant $K(E)=r(\A)$.

The source paper answers the existence part affirmatively when $X$ is compact
metrizable (Theorem~2.12), and gives rank-$3$ counterexamples to the optimal
constant.  Its separate homogeneous-Banach-bundle question is also answered
affirmatively and quantitatively there.  The genuine residue was therefore:
does an expectation always exist when the compact Hausdorff base is allowed
to be nonmetrizable?

\section*{The later counterexample theorem}
For a compact Hausdorff space $Z$ and $n\ge1$, let
\[
 Z_{[n]}=\{S\subseteq Z:1\le |S|\le n\}
\]
with its Vietoris topology, and form the incidence space
\[
 Z^{\subseteq}_{[n]}
   =\{(z,S)\in Z\times Z_{[n]}:z\in S\}.
\]
Projection onto $S$ is an open continuous surjection
\(
 \pi:Z^{\subseteq}_{[n]}\twoheadrightarrow Z_{[n]}
\)
with fibers of size at most $n$, hence a classical branched cover.

Theorem~1.3 of \cite{answer} states that if $n\ge3$ and $Z$ contains the
one-point compactification $\Delta^+$ of an uncountable discrete space
$\Delta$, then the dual unital inclusion
\begin{equation}\label{eq:inclusion}
 C(Z_{[n]})\mathrel{\mathop{\hookrightarrow}^{\pi^*}}
 C(Z^{\subseteq}_{[n]})
\end{equation}
admits \emph{no conditional expectation}.  In bundle language, the fiber at
$S$ is the finite-dimensional abelian algebra $C(S)\cong\C^{|S|}$, so
\eqref{eq:inclusion} is precisely a continuous unital subhomogeneous
$C(Z_{[n]})$-algebra.  It is therefore a counterexample in the exact
generality of Problem~3.11(a), stronger than needed because even an
infinite-index expectation is absent.

\section*{Idea of the counterexample}
An expectation $E:C(Z^{\subseteq}_{[3]})\to C(Z_{[3]})$ determines, at each
finite set $S\in Z_{[3]}$, a state on the fiber $C(S)$.  Identifying points of
$Z$ with their evaluation characters in the weak-$^*$ dual of $C(Z)$, this
is a probability measure
\[
 \mu_S\in\operatorname{conv}(S).
\]
Continuity of $E$ forces $S\mapsto\mu_S$ to be a continuous selection of the
lower-semicontinuous multifunction
\(
 S\mapsto\operatorname{conv}(S)
\).

For $Z=\Delta^+=\Delta\sqcup\{\infty\}$ with $\Delta$ uncountable discrete,
a classical Michael-type selection obstruction rules out such a selection on
the triples $\{\delta,\delta',\infty\}$.  In concrete terms, continuity near
each pair involving $\infty$ forces certain sets of permissible partners to
be cofinite.  An uncountable pigeonhole argument produces one point lying in
infinitely many of those cofinite sets; applying the same continuity
constraint in the reverse direction forces its own set to omit infinitely
many points, a contradiction.  Restricting a hypothetical selection reduces
the cases $Z\supseteq\Delta^+$ and $n\ge3$ to this $Z=\Delta^+$, $n=3$ case.
Consequently no expectation can exist.

\section*{Exact scope after the two papers}
\begin{center}
\small
\begin{tabularx}{\textwidth}{@{}p{0.33\textwidth}X@{}}
\toprule
Issue & Status \\
\midrule
Existence over compact metrizable $X$ & Affirmative in \cite[Theorem~2.12]{source}. \\
Existence over arbitrary compact Hausdorff $X$ & Negative by the incidence example in \cite[Theorem~1.3]{answer}. \\
Optimal equality $K(E)=r(\A)$ & Negative in general already in \cite{source}; counterexamples begin at rank $3$. \\
Homogeneous Banach-bundle unitarizability & Affirmative, quantitatively, in \cite{source}. \\
Extremally disconnected compact bases & Positive with optimal $K=r(\A)$ in \cite[Proposition~1.16]{answer}. \\
\bottomrule
\end{tabularx}
\end{center}

The positive extremally disconnected result is worth separating from the
negative answer: projectivity of such a compact space splits the surjection
from the compact fiber-state space onto the base, yielding a continuous state
choice and hence an expectation.  The answer is thus not simply
``nonmetrizable bases fail''; the obstruction is tied to continuous
selection, and substantial nonmetrizable positive classes remain.

\paragraph{Novelty status.}
This packet records an exact literature resolution, not a new theorem from
the run.  During the upgrade search, the extremally disconnected argument was
independently recovered, but Proposition~1.16 of the same follow-up already
contains it, so no new partial result is claimed.

\paragraph{Human review.}
Check the restated question and metrizable theorem in source PDF pages~2--3;
check the counterexample statement and selection reduction on supporting PDF
page~5, the explicit Problem~3.11 remark on page~6, and the extremally
disconnected result on pages~15--16.  The conclusion depends on the
continuity of the bundle and nonzero finite-dimensional fibers, all of which
hold for the incidence construction.

\enlargethispage{3\baselineskip}
\begin{thebibliography}{9}
\raggedright
\bibitem{source}
A. Chirvasitu,
\emph{Non-commutative branched covers and bundle unitarizability},
arXiv:2409.03531.  See the restatement of Problem~3.11 and Theorem~2.12.

\bibitem{answer}
A. Chirvasitu,
\emph{Finite-index phenomena and the topology of bundle singularities},
arXiv:2409.17807.  See Theorem~1.3, the remarks immediately following it,
and Proposition~1.16.
\end{thebibliography}

\end{document}
